% ===========================================================================
% Preamble: packages, theorem environments, macros.
% Style mirrors the Zinc+ paper and the sister F_2 companion documents
% (sha-f2x-doc, f2-mle-opening-doc).
% ===========================================================================

\usepackage[a4paper,margin=1in]{geometry}
\usepackage[utf8]{inputenc}
\usepackage[T1]{fontenc}
\usepackage{lmodern}
\usepackage{microtype}

\usepackage{amsmath,amssymb,amsthm,mathtools}
\usepackage{booktabs}
\usepackage{array}
\usepackage{tabularx}
\usepackage{enumitem}
\usepackage{xcolor}
\usepackage[colorlinks=true,
            linkcolor={teal!70!black},
            citecolor={teal!70!black},
            urlcolor={teal!70!black}]{hyperref}
\usepackage{thmtools}
\usepackage{cleveref}

% --- Theorem environments --------------------------------------------------
% Every environment shares the `theorem` numbering sequence (one sees
% Definition 5.1, Lemma 5.2, ...) via thmtools' `sibling=theorem`, while
% cleveref still names each reference by its true kind ("Lemma 3.3",
% "Definition 5.1") rather than collapsing them all to "Theorem".
\theoremstyle{plain}
\declaretheorem[name=Theorem, numberwithin=section]{theorem}
\declaretheorem[name=Lemma, sibling=theorem]{lemma}
\declaretheorem[name=Proposition, sibling=theorem]{proposition}
\declaretheorem[name=Corollary, sibling=theorem]{corollary}
\declaretheorem[name=Claim, sibling=theorem]{claim}

\theoremstyle{definition}
\declaretheorem[name=Definition, sibling=theorem]{definition}
\declaretheorem[name=Construction, sibling=theorem]{construction}
\declaretheorem[name=Protocol, sibling=theorem]{protocol}
\declaretheorem[name=Example, sibling=theorem]{example}

\theoremstyle{remark}
\declaretheorem[name=Remark, sibling=theorem]{remark}

% --- Operators -------------------------------------------------------------
\DeclareMathOperator{\MLE}{MLE}
\DeclareMathOperator{\negl}{negl}
\DeclareMathOperator{\eq}{eq}
\DeclareMathOperator{\enc}{Enc}
\DeclareMathOperator{\Merk}{Merkle}
\DeclareMathOperator{\Commit}{Commit}
\DeclareMathOperator{\Open}{Open}
\DeclareMathOperator{\Verify}{Verify}
\DeclareMathOperator{\rate}{rate}
\DeclareMathOperator{\dist}{\Delta}
\DeclareMathOperator{\rep}{Rep}
\DeclareMathOperator{\acc}{Acc}
\DeclareMathOperator{\img}{img}
\DeclareMathOperator{\Hom}{Hom}
\DeclareMathOperator{\spanop}{span}

% --- Number systems --------------------------------------------------------
\newcommand{\F}{\mathbb{F}}
\newcommand{\Ftwo}{\F_2}
\newcommand{\Q}{\mathbb{Q}}
\newcommand{\Z}{\mathbb{Z}}
\newcommand{\N}{\mathbb{N}}
\newcommand{\bP}{\mathbb{P}}

% Binary extension field K = GF(2^128).
\newcommand{\K}{\mathbb{K}}
\newcommand{\Kbits}{128}

% Cell ring R = F_2[X]/(X^D); deployed at D = 32.
\newcommand{\cellring}{R}
\newcommand{\Dval}{D}

% Degree-<D univariate over a coefficient ring S in the cell variable X.
\newcommand{\polyat}[2]{#1[X]_{<#2}}
\newcommand{\Kpoly}{\polyat{\K}{\Dval}}      % K[X]_{<D}
\newcommand{\Ftwopoly}{\polyat{\Ftwo}{\Dval}}% F_2[X]_{<D} = R (as F_2-module)

% Projection psi_alpha : evaluation of the cell variable at alpha in K.
\newcommand{\psia}{\psi_{\alpha}}
\newcommand{\psiz}{\psi_{z}}

% MLE tilde
\newcommand{\mle}[1]{\widetilde{#1}}

% Hadamard product
\newcommand{\had}{\odot}

% Vectors / matrices
\newcommand{\vv}[1]{\mathbf{#1}}
\newcommand{\mat}[1]{\mathsf{#1}}

% Branding for code symbols, usable in text and math mode.
\newcommand{\code}[1]{\ensuremath{\mathsf{#1}}}

\newcommand{\negligible}{\negl}

% Hyperref-friendly section refs
\renewcommand{\sectionautorefname}{Section}
\renewcommand{\subsectionautorefname}{Section}
\renewcommand{\subsubsectionautorefname}{Section}
